% !TEX root = ../main.tex
% !TEX program = lualatex
\documentclass[../main.tex]{subfiles}
\IfSubfilesClassLoaded{\externaldocument{\subfix{../build/main}}}{}
% =============================================================================
% File: 27A_Naval_Justice_and_Standards.tex
% =============================================================================
\begin{document}
\IfSubfilesClassLoaded{\chapter{EDO Study Guide}}{}
\section{Naval Justice and Standards of Conduct}

\subsection{Summary}
\begin{description}
  \item[\textbf{Legal resources.}] Command leaders rely on \ac{quickman}, \acp{sja}, legal officers/clerks, and regional legal support to spot issues early and obtain advice before acting~\autocite{edo-5-1-1-naval-justice-standards-2025}.
  \item[\textbf{Purpose of military law.}] Military law promotes justice, good order and discipline, and national security while reinforcing the obligation to follow the \ac{loac} under stress~\autocite{edo-5-1-1-naval-justice-standards-2025}.
  \item[\textbf{Investigations.}] Commands use preliminary inquiries and command investigations to gather facts, while certain events must be referred to \ac{ncis}~\autocite{edo-5-1-1-naval-justice-standards-2025}.
  \item[\textbf{Discipline tools.}] Non-punitive measures, \ac{njp}, and \ac{adsep} provide graduated responses to misconduct, with defined rights and procedures~\autocite{edo-5-1-1-naval-justice-standards-2025}.
  \item[\textbf{Ethics baseline.}] Standards of conduct, gift rules, political activity limits, and fundraising constraints protect public trust and individual careers~\autocite{edo-5-1-1-naval-justice-standards-2025}.
\end{description}

\subsection{Legal Resources and Roles}
\begin{description}
  \item[\textbf{Staff Judge Advocate.}] The \ac{sja} advises the commander on legal risk, considerations, and options; commanders make the final decision~\autocite{edo-5-1-1-naval-justice-standards-2025}.
  \item[\textbf{Regional legal support.}] \ac{rlso}/\ac{lsss} command services attorneys provide command advice, while \ac{rlso}/\ac{lsss} legal assistance attorneys counsel individuals on civil matters (wills, family law, landlord-tenant, and similar issues)~\autocite{edo-5-1-1-naval-justice-standards-2025}.
  \item[\textbf{Defense and victim counsel.}] The \ac{dso} advises personnel facing disciplinary action, while \ac{vlc} support qualifying victims and coordinate through victim advocates or \ac{sarc} channels~\autocite{edo-5-1-1-naval-justice-standards-2025}.
  \item[\textbf{Command legal enablers.}] Legal officers/clerks and legalmen are unit-level resources, and the \ac{quickman} helps spot legal issues but does not replace \ac{sja} advice~\autocite{edo-5-1-1-naval-justice-standards-2025}.
\end{description}

\subsection{Legal Authorities and Purpose of Military Law}
\begin{description}
  \item[\textbf{Hierarchy of law.}] The Constitution is supreme, Congress enacts statutes, the Executive issues orders and regulations, and military orders derive authority from higher law and leaders (including \ac{secnav} and \ac{cno})~\autocite{edo-5-1-1-naval-justice-standards-2025}.
  \item[\textbf{Purpose.}] Military law promotes justice, deters misconduct, and maintains good order and discipline to strengthen national security and enable \ac{loac} compliance in combat stress~\autocite{edo-5-1-1-naval-justice-standards-2025}.
\end{description}

\subsection{Investigations and Referrals}
\begin{description}
  \item[\textbf{Administrative investigations.}] Common types include preliminary inquiries, command investigations, litigation-report investigations, admiralty letter reports, dual purpose investigations, and courts/boards of inquiry~\autocite{edo-5-1-1-naval-justice-standards-2025}.
  \item[\textbf{Preliminary inquiry.}] A rapid fact-gathering step (typically within three days) led by the convening authority or a designated officer, used to decide on next actions (no action, further investigation, consult a judge advocate, or disciplinary action)~\autocite{edo-5-1-1-naval-justice-standards-2025}.
  \item[\textbf{Command investigation.}] A formal, written investigation led by an investigating officer, usually completed within 30 days, used for significant events (loss/destruction, collisions, some line-of-duty cases, and some deaths) with review by the \ac{gcmca} unless the matter is internal to the unit~\autocite{edo-5-1-1-naval-justice-standards-2025}.
  \item[\textbf{Mandatory \ac{ncis} referral.}] Refer major criminal offenses, sexual assault reports (when not restricted), non-combat deaths not due to natural causes, missing members where foul play is possible, loss/compromise of classified information, and fires or explosions of unknown origin~\autocite{edo-5-1-1-naval-justice-standards-2025}.
\end{description}

\subsection{Rights Advisement and Statements}
\begin{description}
  \item[\textbf{Article 31(b) rights.}] \ac{ucmj} Article 31(b) includes the right to remain silent and advisement on the consequences of waiver; the right to counsel applies to custodial interrogation, and commands should use the standard rights advisement forms~\autocite{edo-5-1-1-naval-justice-standards-2025}.
  \item[\textbf{Cleansing warnings.}] When a prior statement may be inadmissible, a cleansing warning is used; while rules of evidence do not apply to \ac{njp} or \ac{adsep}, rights violations can affect court-martial admissibility and appeals~\autocite{edo-5-1-1-naval-justice-standards-2025}.
\end{description}

\subsection{Search Authorization, Inspections, and Inventories}
\begin{description}
  \item[\textbf{Search basics.}] A search is a government quest for evidence where a reasonable expectation of privacy exists~\autocite{edo-5-1-1-naval-justice-standards-2025}.
  \item[\textbf{Consent v.s. command authorization.}] Best practice is to request consent first (\ac{pass}), with \ac{cass} used when consent is questionable or revoked; obtain authorization in writing when possible~\autocite{edo-5-1-1-naval-justice-standards-2025}.
  \item[\textbf{Probable cause.}] Command-authorized searches require probable cause that a crime occurred and evidence is located in the place to be searched; commanders must not rely solely on another's determination~\autocite{edo-5-1-1-naval-justice-standards-2025}.
  \item[\textbf{Inspections.}] Inspections and inventories are not searches when primarily for military purpose (security, fitness, good order, readiness). Unscheduled or targeted inspections can be improper and risk suppression of evidence~\autocite{edo-5-1-1-naval-justice-standards-2025}.
\end{description}

\subsection{Line of Duty Determinations}
\begin{description}
  \item[\textbf{When required.}] A line-of-duty determination is required for injuries or disease that may result in permanent disability, inability to perform duties for more than 24 hours, or death; a preliminary inquiry is the minimum starting point~\autocite{edo-5-1-1-naval-justice-standards-2025}.
\end{description}

\subsection{Non-Punitive Measures and Non-Judicial Punishment}
\begin{description}
  \item[\textbf{Non-punitive measures.}] Extension of working hours, liberty risk, and extra military instruction are tools for correcting minor misconduct before formal punishment~\autocite{edo-5-1-1-naval-justice-standards-2025}.
  \item[\textbf{NJP authority.}] \acp{co} may award \ac{njp} to all command members (including those on temporary additional duty); \acp{oic} may award \ac{njp} to enlisted members only, and authority cannot be directed but may be withheld for specific offenses or categories~\autocite{edo-5-1-1-naval-justice-standards-2025}.
  \item[\textbf{DRB and XOI.}] The \ac{drb} and \ac{xoi} provide fact-finding and corrective-action screening before \ac{njp}~\autocite{edo-5-1-1-naval-justice-standards-2025}.
  \item[\textbf{Right to refuse.}] Sailors may refuse \ac{njp} unless attached to or embarked on an operational vessel; the right can be exercised until punishment is awarded and does not create a demand for court-martial~\autocite{edo-5-1-1-naval-justice-standards-2025}.
  \item[\textbf{Rights at \ac{njp}.}] The accused may be present, remain silent, have a personal representative, examine evidence, present extenuation/mitigation, and request an open hearing (public is not the same as all-hands mast)~\autocite{edo-5-1-1-naval-justice-standards-2025}.
  \item[\textbf{Appeals.}] \ac{njp} appeals are normally submitted within five working days; members have a right to consult defense counsel, and appeals are limited to unjust or disproportionate punishment reviewed by the \ac{gcmca} for Navy cases~\autocite{edo-5-1-1-naval-justice-standards-2025}.
\end{description}

\subsection{Administrative Separation}
\begin{description}
  \item[\textbf{Purpose.}] \ac{adsep} separates enlisted members who do not meet performance, conduct, or discipline standards after reasonable efforts at counseling, retraining, and rehabilitation~\autocite{edo-5-1-1-naval-justice-standards-2025}.
  \item[\textbf{Reasons.}] Bases include misconduct, unsatisfactory performance, drug and alcohol rehabilitation failure, defective enlistments, security concerns, convenience of the government, and separation in lieu of trial by court-martial~\autocite{edo-5-1-1-naval-justice-standards-2025}.
  \item[\textbf{Common misconduct bases.}] Drug abuse, commission of a serious offense, and a pattern of misconduct (two or more incidents in the same enlistment) are common bases for processing~\autocite{edo-5-1-1-naval-justice-standards-2025}.
  \item[\textbf{Mandatory processing.}] Mandatory bases include sexual misconduct, certain sexual harassment, violent misconduct with severe outcomes, drug paraphernalia or abuse, extremist conduct, family advocacy program failure, alcohol rehabilitation failure, second DUI, and nonconsensual distribution of intimate images~\autocite{edo-5-1-1-naval-justice-standards-2025}.
  \item[\textbf{Procedures and boards.}] Notification procedures apply when the least favorable characterization is honorable or general; board procedures apply when the least favorable characterization is other than honorable and include a formal board unless waived~\autocite{edo-5-1-1-naval-justice-standards-2025}.
  \item[\textbf{Board findings.}] Boards determine whether a basis is met, whether to separate or retain, and the characterization (honorable, general, or other than honorable)~\autocite{edo-5-1-1-naval-justice-standards-2025}.
\end{description}

\subsection{Ethics and Standards of Conduct}
\begin{description}
  \item[\textbf{Ethics foundation.}] The standards of ethical conduct (5 C.F.R. Part 2635) and the \ac{jer} are punitive and enforce the principle that public service is a public trust~\autocite{edo-5-1-1-naval-justice-standards-2025}.
  \item[\textbf{Gifts.}] Employees may not solicit or accept gifts from prohibited sources or because of official position; prohibited sources include entities doing or seeking business with the Department or whose interests are affected by official duties~\autocite{edo-5-1-1-naval-justice-standards-2025}.
  \item[\textbf{Gift exceptions.}] Exceptions include the 20/50 rule, personal relationships, awards, spouse employment, informational materials, and generally available discounts; low-value items like coffee, plaques, or contests are not gifts~\autocite{edo-5-1-1-naval-justice-standards-2025}.
  \item[\textbf{Gifts between employees.}] Subordinates generally should not give gifts to superiors, and employees may not accept gifts from subordinates, with exceptions for modest occasional gifts, shared refreshments, personal hospitality, and special infrequent occasions (marriage, condolence, birth/adoption, or termination of the relationship)~\autocite{edo-5-1-1-naval-justice-standards-2025}.
  \item[\textbf{Group gifts.}] Group gifts to superiors require voluntary contributions, no more than \$10 per contributor, and a total value cap of \$480 (raised from \$300 on 30 Mar 2023)~\autocite{edo-5-1-1-naval-justice-standards-2025}.
\end{description}

\subsection{Fundraising and Political Activity}
\begin{description}
  \item[\textbf{Commercial dealings.}] Personnel should not sell goods to subordinates or their families except in limited personal or off-duty retail contexts without coercion~\autocite{edo-5-1-1-naval-justice-standards-2025}.
  \item[\textbf{Fundraising.}] Do not use uniform, title, or authority to fundraise or fundraise in the workplace, with limited exceptions for \ac{cfc}, \ac{nmcrs}, and similar official campaigns; consult legal counsel before fundraising~\autocite{edo-5-1-1-naval-justice-standards-2025}.
  \item[\textbf{Political activity.}] Members may vote, register, donate, attend rallies as spectators, and express personal opinions, but may not campaign, fundraise, hold public office, distribute partisan material, post political signs in government housing, or wear uniforms to political events~\autocite{edo-5-1-1-naval-justice-standards-2025}.
\end{description}

\subsection{Operational Law}
\begin{description}
  \item[\textbf{Operational discipline.}] The purpose of military law and training is to strengthen readiness and enable compliance with the \ac{loac} during operations~\autocite{edo-5-1-1-naval-justice-standards-2025}.
\end{description}

\subsection{Quick Review}
\begin{itemize}
  \item \ac{quickman} is a spotting tool; \acp{sja} remain the authoritative source of legal advice~\autocite{edo-5-1-1-naval-justice-standards-2025}.
  \item Preliminary inquiries are fast fact-finding tools; command investigations are formal written products for significant events~\autocite{edo-5-1-1-naval-justice-standards-2025}.
  \item \ac{njp} rights include presence, silence, a representative, evidence review, and mitigation opportunities~\autocite{edo-5-1-1-naval-justice-standards-2025}.
  \item \ac{adsep} procedures and characterization thresholds determine whether notification or board processing applies~\autocite{edo-5-1-1-naval-justice-standards-2025}.
  \item Ethics rules prohibit accepting gifts from prohibited sources and restrict political activity and fundraising~\autocite{edo-5-1-1-naval-justice-standards-2025}.
\end{itemize}

\IfSubfilesClassLoaded{
  \RenewDocumentCommand{\entryname}{}{\textbf{\color{Modern} Acronym}}
  \RenewDocumentCommand{\descriptionname}{}{\textbf{\color{Modern} Definition}}
  \printnoidxglossary[
    type=\acronymtype,
    title=Acronyms,
    style=long-booktabs
  ]}{}
\end{document}
